\documentclass[11pt]{article}
\usepackage{mathblatt}
\begin{document}
\blattkopf{Trigonometrie}{Lernblatt}
\weit
\verzeichniszeile{\verz{e1}{1 Seite berechnen mit Sinus, Kosinus und Tangens (Nr.~12–18)} \verztrenn \verz{e2}{2 Winkel berechnen (Nr.~19–22)} \verztrenn \verz{e3}{3 Rechtwinklige Teildreiecke in Figuren und Vermessung (Nr.~23–25)} \verztrenn \verz{e4}{4 Sinussatz (Nr.~26–27)} \verztrenn \verz{abhaken}{Das kann ich}}
% Einheit 1 von 4 – Seite berechnen mit Sinus, Kosinus und Tangens (bank/trigonometrie/e1.jsonl, zusammenbau v0.1)
%% TODO Zweigzeile Teil 1: was hier gelernt wird (Satz in Schülersprache aus dem Einheitstitel) – Einheit 1
\einheitenkopf[e1]{Einheit 1 von 4 · Seite berechnen mit Sinus, Kosinus und Tangens}
\zweigzeile{ab Kl. 10 (am Gymnasium ab Kl. 9) · P10 oft}
\setcounter{aufgabe}{11}

%% TODO Titel: Ich-kann-Satz fehlt in der Bank (Platzhalter: Kettenname) – Nr. 12
%% TODO Anweisung: ein Satz über der ersten Teilaufgabe fehlt in der Bank – Nr. 12
\begin{aufgabe}{sin, cos oder tan?}
\begin{teile}
\teil Rechter Winkel bei $B$. Vom Winkel $\alpha$ aus: $a$ ist gegeben, $b$ ist gesucht. Welche Funktion verbindet $a$ und $b$? \\ \kreuz{sin} \kreuz{cos} \kreuz{tan}

\dreieck{(0,0)}{(5,0)}{(5,3)}{a}{?}{c}{\alpha}{}{}
\end{teile}
\end{aufgabe}

%% TODO Titel: Ich-kann-Satz fehlt in der Bank (Platzhalter: Kettenname) – Nr. 13
%% TODO Anweisung: ein Satz über der ersten Teilaufgabe fehlt in der Bank – Nr. 13
\begin{aufgabe}{Mal oder geteilt?}
\begin{teile}
\teil $\mathrm{sin}\,23^\circ = \frac{x}{8}$, gesucht ist die Seite $x$. Wird mal genommen oder geteilt? \\ \kreuz{im Zähler: mal} \\ \kreuz{im Nenner: geteilt}
\end{teile}
\end{aufgabe}

%% TODO Titel: Ich-kann-Satz fehlt in der Bank (Platzhalter: Kettenname) – Nr. 14
%% TODO Anweisung: ein Satz über der ersten Teilaufgabe fehlt in der Bank – Nr. 14
\begin{aufgabe}{Seite berechnen}
%% TODO \rechenplatz{n}: Schrittzahl des Grundfalls fehlt in der Bank (nur bei fester Schreibform, 2.3 b)
\begin{teile}
\teil Rechter Winkel bei $B$. Schreibe an jede Seite H, G oder A – vom Winkel $\alpha$ aus.

\dreieck{(0,0)}{(5,0)}{(5,3)}{r}{t}{s}{\alpha}{}{}
\teil Rechter Winkel bei $C$. Trage den Bruch ein. \hfill (P10 2017 OS) \\

\dreieck{(0,0)}{(5,0)}{(1.8,2.4)}{d}{e}{f}{}{\beta}{}
$\mathrm{sin}\,\beta =$ \leerfeld
\teil Rechter Winkel bei $A$. Welche Gleichung gilt? \hfill (P10 2018 OS) \\ \kreuz{$\mathrm{sin}\,\beta = \frac{p}{q}$} \\ \kreuz{$\mathrm{cos}\,\beta = \frac{q}{p}$} \\ \kreuz{$\mathrm{sin}\,\beta = \frac{q}{p}$}

\dreieck{(0,0)}{(4,0)}{(0,3)}{p}{q}{r}{}{\beta}{}
\teil Taschenrechner im Grad-Modus: $\mathrm{sin}\,23^\circ$ auf drei Dezimalen? \leerfeld
\teil Hypotenuse $8$ cm, $\alpha = 31^\circ$ – wie lang ist die Gegenkathete $x$ von $\alpha$? x ≈ \leerfeld[cm]
\teil Hypotenuse $9$ cm, $\alpha = 27^\circ$ – wie lang ist die Ankathete $x$ von $\alpha$? x ≈ \leerfeld[cm]
\teil Rechter Winkel bei $A$, $a = 7$ cm, $\beta = 41^\circ$. Wie lang ist $b$?

\dreieck{(0,0)}{(4,0)}{(0,3)}{a}{b}{c}{}{\beta}{}
b ≈ \leerfeld[cm]
\teil Gegenkathete $6$ cm, $\alpha = 44^\circ$ – wie lang ist die Hypotenuse $x$? x ≈ \leerfeld[cm]
\teil Ankathete $7$ cm, $\alpha = 32^\circ$ – wie lang ist die Hypotenuse $x$? x ≈ \leerfeld[cm]
\teil Ankathete $6$ cm, $\alpha = 27^\circ$ – wie lang ist die Gegenkathete $x$? x ≈ \leerfeld[cm]
\teil Gegenkathete $5$ cm, $\alpha = 31^\circ$ – wie lang ist die Ankathete $x$? x ≈ \leerfeld[cm]
\teil Hypotenuse $9$ cm, $\alpha = 72^\circ$ – wie lang ist die Ankathete $x$? Wähle die Funktion selbst. x ≈ \leerfeld[cm]
\end{teile}
\begin{gleichungsraster}[2]
\gl{\text{$\mathrm{sin}\,48^\circ = \frac{6}{x}$ – wie groß ist $x$?}} &
\end{gleichungsraster}
\begin{teile}
\teil Hypotenuse $6{,}8$ m, $\alpha = 42{,}6^\circ$ – wie lang ist die Gegenkathete $x$? x ≈ \leerfeld[m]
\teil Hypotenuse $7$ cm, $\alpha = 42^\circ$. Zeige rechnerisch, dass die Ankathete $x$ rund $5{,}2$ cm lang ist.
\teil Eine $4{,}2$ m lange Leiter lehnt an einer Wand und bildet mit dem Boden einen Winkel von $71^\circ$ (Skizze). In welcher Höhe liegt sie an der Wand an?

\dreieck{(0,0)}{(1.4,0)}{(1.4,4)}{Wand}{Leiter}{Boden}{71^\circ}{}{}
\leerfeld[m]
\teil Ein Grundstück $ABCD$ hat bei $A$ einen rechten Winkel. Die Diagonale $BD$ ist $840$ m lang, der Winkel $ADB$ beträgt $47^\circ$. Zeige rechnerisch, dass $AB \approx 614$ m ist. \hfill (P10 2021 OS)
\end{teile}
\end{aufgabe}

%% TODO Titel: Ich-kann-Satz fehlt in der Bank (Platzhalter: Kettenname) – Nr. 15
%% TODO Anweisung: ein Satz über der ersten Teilaufgabe fehlt in der Bank – Nr. 15
\begin{aufgabe}{Seiten vom zweiten spitzen Winkel aus neu benennen (Gegen- und Ankathete tauschen)}
\begin{teile}
\teil Rechter Winkel bei $C$. Vom Winkel $\alpha$ aus ist $a$ die Gegenkathete und $b$ die Ankathete. Wie heißen $a$ und $b$ vom Winkel $\beta$ aus? \leerfeld und \leerfeld
\end{teile}
\end{aufgabe}

%% TODO Titel: Ich-kann-Satz fehlt in der Bank (Platzhalter: Kettenname) – Nr. 16
%% TODO Anweisung: ein Satz über der ersten Teilaufgabe fehlt in der Bank – Nr. 16
\begin{aufgabe}{Seitenverhältnis aus gegebenen Längen als Dezimalzahl und Vergleich mit dem Taschenrechnerwert des Winkels}
\begin{teile}
\teil Gemessen: Gegenkathete $3{,}4$ cm, Hypotenuse $5$ cm, Winkel $43^\circ$. Verhältnis als Dezimalzahl – passt es zu $\mathrm{sin}\,43^\circ$? \leerfeld
\end{teile}
\end{aufgabe}

%% TODO Titel: Ich-kann-Satz fehlt in der Bank (Platzhalter: Kettenname) – Nr. 17
%% TODO Anweisung: ein Satz über der ersten Teilaufgabe fehlt in der Bank – Nr. 17
\begin{aufgabe}{Ergebnis prüfen: Kathete kürzer als Hypotenuse}
\begin{teile}
\teil Lea rechnet zur Hypotenuse $6$ cm eine Gegenkathete von $7{,}3$ cm aus. Kann das stimmen?
\end{teile}
\end{aufgabe}

%% TODO Titel: Ich-kann-Satz fehlt in der Bank (Platzhalter: Kettenname) – Nr. 18
%% TODO Anweisung: ein Satz über der ersten Teilaufgabe fehlt in der Bank – Nr. 18
\begin{aufgabe}{Seite berechnen – Fehler finden, Begründen, Anwendung}
\begin{teile}
\teil Gegenkathete $5$ cm, $\alpha = 41^\circ$, gesucht die Hypotenuse $x$. Wo ist der Fehler? \rechnung{\mathrm{sin}\,41^\circ &= \frac{5}{x} \\ x &= 5 \cdot \mathrm{sin}\,41^\circ \\ x &\approx 3{,}3 \text{ cm}}
\teil Warum ist $\mathrm{sin}\,\alpha$ im rechtwinkligen Dreieck immer kleiner als $1$?
\teil Eine $6$ m lange Leiter soll mit $72^\circ$ zum Boden an der Hauswand stehen. Das Fenster beginnt in $5{,}5$ m Höhe. Reicht die Leiter bis zum Fenster?
\end{teile}
\end{aufgabe}

\clearpage
% Einheit 2 von 4 – Winkel berechnen (bank/trigonometrie/e2.jsonl, zusammenbau v0.1)
%% TODO Zweigzeile Teil 1: was hier gelernt wird (Satz in Schülersprache aus dem Einheitstitel) – Einheit 2
\einheitenkopf[e2]{Einheit 2 von 4 · Winkel berechnen}
\zweigzeile{ab Kl. 10 (am Gymnasium ab Kl. 9) · P10 oft}
\setcounter{aufgabe}{18}

%% TODO Titel: Ich-kann-Satz fehlt in der Bank (Platzhalter: Kettenname) – Nr. 19
%% TODO Anweisung: ein Satz über der ersten Teilaufgabe fehlt in der Bank – Nr. 19
\begin{aufgabe}{Winkel berechnen}
%% TODO \rechenplatz{n}: Schrittzahl des Grundfalls fehlt in der Bank (nur bei fester Schreibform, 2.3 b)
\begin{teile}
\teil Gegeben: Hypotenuse $7$ cm und Gegenkathete $4$ cm. Gesucht: der Winkel $\alpha$. Was ist gesucht? \\ \kreuz{Seite: Gleichung umstellen} \\ \kreuz{Winkel: Umkehrtaste}
\teil Gegenkathete $3$ cm, Hypotenuse $8$ cm – wie groß ist $\alpha$? $\alpha \approx$ \leerfeld $^\circ$
\teil Ankathete $5$ cm, Hypotenuse $7$ cm – wie groß ist $\alpha$? $\alpha \approx$ \leerfeld $^\circ$
\teil Gegenkathete $5$ cm, Ankathete $8$ cm – wie groß ist $\alpha$? $\alpha \approx$ \leerfeld $^\circ$
\teil Rechter Winkel bei $A$, $a = 10$ cm, $b = 3$ cm. Wie groß ist $\beta$?

\dreieck{(0,0)}{(4,0)}{(0,3)}{a}{b}{c}{}{\beta}{}
$\beta \approx$ \leerfeld $^\circ$
\teil Gegenkathete von $\alpha$: $3$ cm, Hypotenuse $5$ cm. Berechne $\alpha$, dann $\beta$ als Ergänzung; kontrolliere $\beta$ mit dem Kosinus. $\alpha \approx$ \leerfeld $^\circ$, $\beta \approx$ \leerfeld $^\circ$
\teil Gegenkathete $2{,}3$ m, Hypotenuse $4{,}7$ m – wie groß ist $\alpha$, auf eine Dezimale? $\alpha \approx$ \leerfeld $^\circ$
\teil Eine Seilbahn überwindet $420$ m Höhe, das Seil ist $1\,300$ m lang. Wie groß ist der Steigungswinkel? \leerfeld $^\circ$
\teil Die Punkte $A$, $B$, $C$ bilden ein Dreieck mit rechtem Winkel bei $A$. Lies die Längen der Katheten ab und berechne den Winkel $\beta$ bei $B$.

\begin{ksys}[xmin=-1,xmax=6,ymin=-1,ymax=5,ablesen] \punkt{1}{1}{A} \punkt{5}{1}{B} \punkt{1}{4}{C} \end{ksys}
$\beta \approx$ \leerfeld $^\circ$
\teil Im Dreieck $ABC$ ist die Höhe von $C$ auf $AB$ $4{,}5$ cm lang, $AC = 7{,}2$ cm. Wie groß ist $\alpha$ bei $A$? $\alpha \approx$ \leerfeld $^\circ$
\teil Ankathete $5{,}2$ m, Hypotenuse $8$ m. Zeige rechnerisch, dass $\alpha$ rund $49^\circ$ groß ist.
\teil Eine Seilbahn führt von der Talstation $A$ zur Bergstation $B$. Das Seil ist $610$ m lang, die waagerechte Strecke unter dem Seil $460$ m. Berechne den Winkel zwischen Seil und Waagerechter. \hfill (P10 2024 OS) \\ \leerfeld $^\circ$
\end{teile}
\end{aufgabe}

%% TODO Titel: Ich-kann-Satz fehlt in der Bank (Platzhalter: Kettenname) – Nr. 20
%% TODO Anweisung: ein Satz über der ersten Teilaufgabe fehlt in der Bank – Nr. 20
\begin{aufgabe}{Winkel prüfen: die längere Kathete liegt dem größeren Winkel gegenüber}
\begin{teile}
\teil Die Katheten sind $4$ cm und $9$ cm lang. Tim sagt: Der Winkel gegenüber der $4$-cm-Kathete ist $66^\circ$. Kann das stimmen?
\end{teile}
\end{aufgabe}

%% TODO Titel: Ich-kann-Satz fehlt in der Bank (Platzhalter: Kettenname) – Nr. 21
%% TODO Anweisung: ein Satz über der ersten Teilaufgabe fehlt in der Bank – Nr. 21
\begin{aufgabe}{Umkehrtaste am Taschenrechner (SHIFT sin, cos, tan; Anzeige sin⁻¹)}
\begin{teile}
\teil Mit der Umkehrtaste: $\mathrm{sin}^{-1}(0{,}42)$ – welcher Winkel, auf eine Dezimale? \leerfeld $^\circ$
\end{teile}
\end{aufgabe}

%% TODO Titel: Ich-kann-Satz fehlt in der Bank (Platzhalter: Kettenname) – Nr. 22
%% TODO Anweisung: ein Satz über der ersten Teilaufgabe fehlt in der Bank – Nr. 22
\begin{aufgabe}{Winkel berechnen – Fehler finden, Begründen, Anwendung}
\begin{teile}
\teil Gegenkathete von $\alpha$: $3$ cm, Ankathete: $8$ cm. Wo ist der Fehler? \rechnung{\mathrm{tan}\,\alpha &= \frac{8}{3} \\ \alpha &\approx 69{,}4^\circ}
\teil Ole tippt $\mathrm{sin}^{-1}(1{,}3)$ und der Taschenrechner meldet einen Fehler. Warum?
\teil Eine Rollstuhlrampe darf höchstens $3{,}4^\circ$ steil sein. Eine $5$ m lange Rampe überwindet $0{,}32$ m. Ist sie zulässig?
\end{teile}
\end{aufgabe}

\clearpage
% Einheit 3 von 4 – Rechtwinklige Teildreiecke in Figuren und Vermessung (bank/trigonometrie/e3.jsonl, zusammenbau v0.1)
%% TODO Zweigzeile Teil 1: was hier gelernt wird (Satz in Schülersprache aus dem Einheitstitel) – Einheit 3
\einheitenkopf[e3]{Einheit 3 von 4 · Rechtwinklige Teildreiecke in Figuren und Vermessung}
\zweigzeile{ab Kl. 10 (am Gymnasium ab Kl. 9) · P10 oft}
\setcounter{aufgabe}{22}

%% TODO Titel: Ich-kann-Satz fehlt in der Bank (Platzhalter: Kettenname) – Nr. 23
%% TODO Anweisung: ein Satz über der ersten Teilaufgabe fehlt in der Bank – Nr. 23
\begin{aufgabe}{Teildreiecke und Vermessung}
%% TODO \rechenplatz{n}: Schrittzahl des Grundfalls fehlt in der Bank (nur bei fester Schreibform, 2.3 b)
\begin{teile}
\teil Fahre im Trapez das rechtwinklige Teildreieck mit der Höhe $h$ links nach. Beschrifte seine Seiten mit H, G, A – vom Basiswinkel unten links aus.

\trapez[hoehe=h]{5}{3}{2.5}
\teil Gleichschenkliges Dreieck: Schenkel $7$ cm, Basiswinkel $63^\circ$. Wie lang ist die Höhe $h$ auf die Grundseite? h ≈ \leerfeld[cm]
\teil Trapez: Höhe $4$ m, Basiswinkel $62^\circ$. Wie lang ist der Schenkel $s$? s ≈ \leerfeld[m]
\teil Trapez: Schenkel $5$ cm, Basiswinkel $58^\circ$. Wie hoch ist das Trapez? h ≈ \leerfeld[cm]
\teil Rechtwinkliges Trapez: die senkrechten Seiten sind $9$ cm und $4$ cm lang und stehen $8$ cm auseinander. Wie groß ist der Winkel $\alpha$ zwischen der schrägen Seite und der Waagerechten? $\alpha \approx$ \leerfeld $^\circ$
\teil Parallelogramm $ABCD$: $BC = 9$ cm. $F$ liegt auf der Verlängerung von $AB$ über $B$ hinaus, $CF$ steht senkrecht darauf; der Winkel $CBF$ beträgt $58^\circ$. Wie lang ist die Höhe $h = CF$? \leerfeld[cm]
\teil Parallelogramm $ABCD$ mit der Höhe $h = 6$ cm von $C$ auf die Verlängerung von $AB$. Die Diagonale $AC$ bildet mit $AB$ einen Winkel von $27^\circ$. Wie lang ist $AC$? AC ≈ \leerfeld[cm]
\teil Drachen: eine obere Seite ist $6$ cm lang und bildet mit der Längsdiagonale einen Winkel von $32^\circ$. Wie lang ist die halbe Querdiagonale? \leerfeld[cm]
\teil Aus $45$ m waagerechtem Abstand sieht man die Spitze eines Baums unter einem Höhenwinkel von $31^\circ$. Wie hoch liegt sie über dem Messgerät? \leerfeld[m]
\teil Ein Messgerät steht in $1{,}4$ m Höhe, $58$ m waagerecht entfernt. Der Höhenwinkel zur Spitze eines Kirchturms beträgt $27^\circ$. Wie hoch ist es? \leerfeld[m]
\teil Aus $48$ m Abstand peilt man in $1{,}7$ m Augenhöhe die Dachkante eines Turms unter $24^\circ$ an. Auf dem Dach steht eine $6$ m hohe Antenne. Wie hoch liegt die Antennenspitze über dem Boden? \leerfeld[m]
\teil Rechter Winkel bei $D$; $B$ liegt zwischen $D$ und $C$. Von $A$ aus: $\angle DAC = 71^\circ$, $\angle BAC = 23^\circ$, $AD = 32$ m. Wie lang ist $BD$? BD ≈ \leerfeld[m]
\teil Dreieck $ABC$ mit der Höhe $h$ von $C$ auf $AB$: $AC = 8$ cm, $\alpha = 42^\circ$, $\beta = 61^\circ$. Berechne erst $h$, dann $BC$. h ≈ \leerfeld[cm] , BC ≈ \leerfeld[cm]
\teil Gleichschenkliges Trapez: bekannt sind die beiden parallelen Seiten und die Höhe, gesucht ist der Umfang. Beschreibe den Rechenweg ohne zu rechnen.
\teil Hypotenuse $9$ cm, $\alpha = 33^\circ$. Berechne Gegenkathete und Ankathete und prüfe mit dem Satz des Pythagoras. G ≈ \leerfeld[cm] , A ≈ \leerfeld[cm]
\teil Dreieck $ABC$: $c = 9$ cm, $b = 6$ cm, $\alpha = 47^\circ$. Berechne die Höhe $h$ auf $c$ und dann den Flächeninhalt. h ≈ \leerfeld[cm] , A ≈ \leerfeld[cm²]
\teil Quadratische Pyramide: Grundkante $6$ m, die Seitenflächen sind um $57^\circ$ gegen die Grundfläche geneigt. Wie lang ist die Seitenhöhe $h_s$? \leerfeld[m]
\teil Ein Messgerät steht $1{,}4$ m über dem Boden. Die Sichtlinie zur Hügelkuppe ist $96{,}5$ m lang und steigt unter $27^\circ$ an. Wie hoch ist der Hügel? Auf der Kuppe steht ein $9$ m hoher Aussichtsplattform – wie hoch über dem Boden liegt seine Oberkante? \hfill (P10 2018 OS) \\ \leerfeld[m] ; \leerfeld[m]
\end{teile}
\end{aufgabe}

%% TODO Titel: Ich-kann-Satz fehlt in der Bank (Platzhalter: Kettenname) – Nr. 24
%% TODO Anweisung: ein Satz über der ersten Teilaufgabe fehlt in der Bank – Nr. 24
\begin{aufgabe}{Strecke aus Teilstrecken nach der Trigonometrie}
\begin{teile}
\teil Dreieck $ABC$ mit der Höhe $CD = 4$ cm auf $AB$: $AD = 6$ cm, $\beta = 53^\circ$. Wie lang ist $AB$? AB ≈ \leerfeld[cm]
\end{teile}
\end{aufgabe}

%% TODO Titel: Ich-kann-Satz fehlt in der Bank (Platzhalter: Kettenname) – Nr. 25
%% TODO Anweisung: ein Satz über der ersten Teilaufgabe fehlt in der Bank – Nr. 25
\begin{aufgabe}{Teildreiecke und Vermessung – Fehler finden, Begründen, Anwendung}
\begin{teile}
\teil Abstand zum Turm $44$ m, Höhenwinkel $31^\circ$, das Messgerät steht $1{,}7$ m hoch. Wo ist der Fehler? \rechnung{h &= 44 \cdot \mathrm{tan}\,31^\circ \\ h &\approx 26{,}4 \text{ m}}
\teil Warum rechnet man im gleichschenkligen Dreieck mit der halben Grundseite, wenn man den Basiswinkel benutzt?
\teil Ein Baum soll gefällt werden, das Haus steht $18$ m vom Stamm entfernt. Aus $20$ m Abstand sieht man die Baumspitze in $1{,}7$ m Augenhöhe unter $41^\circ$. Kann der Baum beim Fallen das Haus treffen?
\end{teile}
\end{aufgabe}

\clearpage
% Einheit 4 von 4 – Sinussatz (bank/trigonometrie/e4.jsonl, zusammenbau v0.1)
%% TODO Zweigzeile Teil 1: was hier gelernt wird (Satz in Schülersprache aus dem Einheitstitel) – Einheit 4
\einheitenkopf[e4]{Einheit 4 von 4 · Sinussatz}
\zweigzeile{ab Kl. 10 (am Gymnasium ab Kl. 9) · P10 oft}
\setcounter{aufgabe}{25}

%% TODO Titel: Ich-kann-Satz fehlt in der Bank (Platzhalter: Kettenname) – Nr. 26
%% TODO Anweisung: ein Satz über der ersten Teilaufgabe fehlt in der Bank – Nr. 26
\begin{aufgabe}{Sinussatz}
%% TODO \rechenplatz{n}: Schrittzahl des Grundfalls fehlt in der Bank (nur bei fester Schreibform, 2.3 b)
\begin{teile}
\teil Im Dreieck $ABC$ ist bei $B$ ein rechter Winkel markiert, $a$ und $\alpha$ sind gegeben. Womit rechnest du $b$? \\ \kreuz{sin, cos, tan} \kreuz{Sinussatz}
\teil Dreieck $ABC$: $a = 7$ cm, $\alpha = 68^\circ$, $\beta = 51^\circ$. Wie lang ist $b$? b ≈ \leerfeld[cm]
\teil Dreieck $PQR$: $PQ = 8$ cm, Winkel bei $R$: $69^\circ$, Winkel bei $Q$: $53^\circ$. Wie lang ist $PR$? PR ≈ \leerfeld[cm]
\teil Dreieck $ABC$: $a = 8$ cm, $\alpha = 63^\circ$, $\beta = 48^\circ$. Wie lang ist $c$? c ≈ \leerfeld[cm]
\teil Im Drachen $ABCD$ gilt im Dreieck $ABD$: $AD = 4{,}6$ cm, der Winkel bei $A$ ist $118^\circ$, der Winkel bei $B$ ist $33^\circ$. Berechne die Diagonale $BD$. \hfill (P10 2015 OS) \\ BD ≈ \leerfeld[cm]
\teil Dreieck $ABC$: $a = 7$ cm, $\alpha = 27^\circ$, $\beta = 118^\circ$. Wie lang ist $c$? c ≈ \leerfeld[cm]
\teil Dreieck $ABC$: $a = 4{,}35$ m, $\alpha = 58{,}2^\circ$, $\beta = 47{,}6^\circ$. Wie lang ist $b$, auf eine Dezimale? b ≈ \leerfeld[m]
\teil Dreieck $ABC$: $a = 9$ cm, $\alpha = 77^\circ$, $\beta = 44^\circ$. Zeige rechnerisch, dass $b \approx 6{,}4$ cm ist.
\teil Dreieck $ABC$: $c = 10$ cm, $\alpha = 53^\circ$, $\beta = 56^\circ$. Prüfe die Aussage „$a$ ist länger als $b$.“
\teil Viereck $ABCD$ mit rechtem Winkel bei $B$ und der Diagonale $AC = 8$ m. $\angle BAC = 57^\circ$, $\angle BCD = 91^\circ$, $\angle ADC = 76^\circ$. Wie lang ist $AD$? AD ≈ \leerfeld[m]
\teil Wanderung: Hütte $H$, Gipfel $G$, See $S$. $HG = 3\,200$ m, der Winkel bei $H$ ist $57^\circ$, der Winkel bei $S$ ist $72^\circ$. Wie lang ist der Weg vom Gipfel zum See und weiter zum Parkplatz, der $1\,800$ m hinter dem See liegt? \hfill (P10 2014 OS) \\ \leerfeld[km]
\teil Dreieck $ABC$: $a = 8$ cm, $\alpha = 67^\circ$, $\beta = 53^\circ$. Berechne $b$ mit dem Sinussatz und noch einmal über die Höhe $h$ von $C$ auf $AB$. b ≈ \leerfeld[cm]
\teil Dreieck $ABC$: $a = 9$ cm, $\alpha = 68^\circ$, $b = 7$ cm. Wie groß ist $\beta$? $\beta \approx$ \leerfeld $^\circ$
\teil Dreieck $ABC$: $a = 5$ cm, $b = 8$ cm, $\gamma = 47^\circ$. Wie lang ist $c$? \leerfeld[cm]
\teil Dreieck $ABC$: $a = 5$ cm, $b = 7$ cm, $c = 6$ cm. Wie groß ist $\gamma$? $\gamma \approx$ \leerfeld $^\circ$
\teil Die Höhe $h_c$ zerlegt das Dreieck $ABC$ in zwei rechtwinklige Teildreiecke. Drücke $h_c$ in beiden Teildreiecken aus und leite daraus $\frac{a}{\mathrm{sin}\,\alpha} = \frac{b}{\mathrm{sin}\,\beta}$ her.
\teil Eine Seilbahn führt von $A$ über die Stütze $B$ nach $C$. $AB = 520$ m, der Winkel bei $A$ ist $31^\circ$, der Winkel $ABC$ ist $113^\circ$. Berechne die Länge der Strecke $BC$. \hfill (P10 2024 OS) \\ BC ≈ \leerfeld[m]
\end{teile}
\end{aufgabe}

%% TODO Titel: Ich-kann-Satz fehlt in der Bank (Platzhalter: Kettenname) – Nr. 27
%% TODO Anweisung: ein Satz über der ersten Teilaufgabe fehlt in der Bank – Nr. 27
\begin{aufgabe}{Sinussatz – Fehler finden, Begründen, Anwendung}
\begin{teile}
\teil Dreieck $ABC$: $a = 8$ cm, $\alpha = 71^\circ$, $\beta = 43^\circ$, gesucht $b$. Wo ist der Fehler? \rechnung{b &= 8 \cdot \mathrm{sin}\,71^\circ : \mathrm{sin}\,43^\circ \\ b &\approx 11{,}1 \text{ cm}}
\teil Die Höhe $h$ von $C$ auf $AB$ ist in beiden Teildreiecken dieselbe Kathete. Warum folgt daraus $b \cdot \mathrm{sin}\,\alpha = a \cdot \mathrm{sin}\,\beta$?
\teil Über einen Fluss soll ein Seil von $A$ zu einem Baum $C$ am anderen Ufer gespannt werden. Am eigenen Ufer ist $AB = 80$ m abgesteckt, der Winkel bei $A$ ist $72^\circ$, der bei $B$ $63^\circ$. Reicht eine $100$ m lange Seilrolle?
\end{teile}
\end{aufgabe}

% Abhakseite „Das kann ich“ – Zone nur im Gesamt (\mitzone)
%% TODO Ich-kann-Sätze: Platzhalter wie in den Aufgabendateien
\begin{abhakseite}
\ifdefined\mitzone
\abhakgruppe{Kennst du schon}
\abhak{1}{Rechtwinkliges Dreieck erkennen, Rechtwinkelmarke (Bogen mit Punkt), Hypotenuse gegenüber dem rechten Winkel, Katheten am rechten Winkel}
\abhak{2}{Winkel benennen (α, β, γ, auch β1, β2, ε), Winkel messen und schätzen, spitz oder stumpf}
\abhak{3}{Taschenrechner}
\abhak{4}{Gleichung mit einem Bruch nach einer Größe umstellen, Größe im Zähler (mal) oder im Nenner (geteilt), mit Strich}
\abhak{5}{Bruch als Verhältnis lesen und in eine Dezimalzahl umwandeln}
\abhak{6}{Satz des Pythagoras als zweiter Weg zur dritten Seite, Vergleich der Ergebnisse}
\abhak{7}{Winkelsumme im Dreieck, die beiden spitzen Winkel des rechtwinkligen Dreiecks als Ergänzung, Winkel aus Teilwinkeln (Differenz), Strecke aus Teilstrecken}
\abhak{8}{Längeneinheiten vor dem Rechnen angleichen (m und km, cm und m), Ergebnis mit Einheit, Runden auf eine Dezimale}
\abhak{9}{Höhe, Diagonale, Überstand am Trapez, Fußpunkt auf der Verlängerung beim Parallelogramm, halbe Diagonale im Drachen}
\abhak{10}{Gleichung mit einem Bruch nach einer Größe umstellen, Größe im Zähler (mal) oder im Nenner (geteilt), mit Strich – Fehler finden}
\abhak{11}{Gleichung mit einem Bruch nach einer Größe umstellen, Größe im Zähler (mal) oder im Nenner (geteilt), mit Strich – selbst rechnen}
\fi
\abhakgruppe{Einheit 1 · Seite berechnen mit Sinus, Kosinus und Tangens}
\abhak{12}{sin, cos oder tan?}
\abhak{13}{Mal oder geteilt?}
\abhak{14}{Seite berechnen}
\abhak{15}{Seiten vom zweiten spitzen Winkel aus neu benennen (Gegen- und Ankathete tauschen)}
\abhak{16}{Seitenverhältnis aus gegebenen Längen als Dezimalzahl und Vergleich mit dem Taschenrechnerwert des Winkels}
\abhak{17}{Ergebnis prüfen: Kathete kürzer als Hypotenuse}
\abhak{18}{Seite berechnen – Fehler finden, Begründen, Anwendung}
\abhakgruppe{Einheit 2 · Winkel berechnen}
\abhak{19}{Winkel berechnen}
\abhak{20}{Winkel prüfen: die längere Kathete liegt dem größeren Winkel gegenüber}
\abhak{21}{Umkehrtaste am Taschenrechner (SHIFT sin, cos, tan; Anzeige sin⁻¹)}
\abhak{22}{Winkel berechnen – Fehler finden, Begründen, Anwendung}
\abhakgruppe{Einheit 3 · Rechtwinklige Teildreiecke in Figuren und Vermessung}
\abhak{23}{Teildreiecke und Vermessung}
\abhak{24}{Strecke aus Teilstrecken nach der Trigonometrie}
\abhak{25}{Teildreiecke und Vermessung – Fehler finden, Begründen, Anwendung}
\abhakgruppe{Einheit 4 · Sinussatz}
\abhak{26}{Sinussatz}
\abhak{27}{Sinussatz – Fehler finden, Begründen, Anwendung}
\end{abhakseite}

\end{document}
